\subsection{Biologically-Inspired AI: Gleanings from Neuroscience and Cognitive Psychology}
\label{sec:4.1.2}
Backpropagation, the mechanism that trains most modern neural networks, is a distant cousin of synaptic weight adjustment in the brain: both strengthen and weaken connections in response to error. That distant-cousin relationship runs through everything below, and an AI system built on a wrong model of how a mind works will fail in ways that are hard to predict from the wrong model alone.

\runin{Where a borrowing stopped}
One borrowing did not survive scrutiny at all. Mirror neurons — cells found in macaque premotor cortex in the early 1990s that fire both when the animal grasps an object and when it watches someone else do it \autocite{gallese1996action} — were cast as the neural basis of empathy, language, imitation and cultural transmission at once \autocite{ramachandran2009neurons}, and the interpretation outran the finding \autocite{hickok2014myth}. Single-neuron evidence was never obtained in humans on the same terms \autocite{mukamel2010singleneuron}, and knowing what someone is doing and knowing what they believe are different computations the brain appears to run in different places. What survives is more modest and still useful: a system that maps an observed action onto its own action repertoire supports imitation learning and gesture recognition, with no claim about what it is like to have one. A single vivid mechanism, well named, absorbed four distinct capacities that later had to be separated again, and architectures for artificial empathy are exposed to the same failure.

\runin{Attention decides what is seen, which makes its failures blind spots}
About half the people asked to count basketball passes on a video will not notice a person in a gorilla costume walk into frame and thump their chest for nine seconds \autocite{simons1999gorillas}. The effect is inattentional blindness, and it does not spare expertise: radiologists searching lung images for nodules miss a gorilla deliberately embedded in the scan 83 percent of the time \autocite{drew2013invisible}.

The design consequence turns on how attention is modeled. If attention is a resource allocator bolted onto a system that has already seen everything, its failures are inefficiencies. If attention is the selection that decides what gets perceived at all, its failures are blind spots — and a system that has learned which regions of its input reward processing has also learned which regions to be blind to. The human case establishes the property that makes this dangerous: the blind spot is invisible from inside, and the viewers who missed the gorilla did not experience a gap. Self-attention has the same shape \autocite{vaswani2017attention}. A model that weights some tokens heavily weights others near zero, nothing in the architecture announces which relevant thing was assigned a low weight, and a system meant to act ethically has to treat the regions it has learned to discount as an always-open question.

\runin{Perception is inference, so the priors are the exposure}
Predictive coding makes the point precise: higher, context-rich areas send predictions down to lower, sensory ones, and the lower areas return only what the predictions failed to anticipate \autocite{rao1999predictive}. What propagates upward is the surprise, and the anatomy fits, backward connections outnumbering forward ones. Anil Seth's phrase for the result is controlled hallucination: the difference between seeing a room and hallucinating one is whether corrections are arriving and being applied \autocite{seth2021being}. A system's perception is only as reliable as the priors it perceives through, so priors trained on skewed or incomplete data will hallucinate in the specific direction of that skew, confidently, with no internal signal that anything is wrong. That is why bias in a perceptual system cannot be caught by asking the system.

\runin{Memory is rebuilt}
Memory is not played back. Frederic Bartlett established this in the 1930s by having people retell an unfamiliar folk tale over weeks: they did not lose pieces the way a recording degrades, they remodeled the story, replacing strange details with ordinary ones \autocite{bartlett1932remembering}. Elizabeth Loftus gave the finding courtroom consequences, showing that the verb used to ask about a filmed car accident changes the speed people report and, a week later, whether they remember broken glass that was never in the film \autocite{loftus1974reconstruction}. Later work implanted whole episodes that never happened. The subjects were not lying, and none of it correlates with felt confidence.

This book asks what it would take to trust what a machine witness reports, and answers that no witness has an honest memory in the strong sense. The design point is that retrieval has to be built expecting reconstruction, and a system that treats a retrieved memory as ground truth inherits the overconfidence of Loftus's subjects with none of the excuse.

\runin{Control without a controller}
The textbook division of labor across prefrontal regions is tempting to read as a committee: several named areas, each handling one job, together adding up to executive control. That reading has a name — the homunculus — and a specific failure mode. If some region is a controller that decides what to attend to, inhibit, and hold in mind, then intelligent coordination has been explained by positing an intelligent coordinator inside the system, which explains nothing: the controller would itself need to attend, decide, and remember, and would need its own controller, without end.

Earl Miller and Jonathan Cohen proposed a model that gets control without a controller \autocite{miller2001integrative}. The prefrontal cortex does not house an agent issuing commands. It maintains representations of goals — sustained patterns of activity encoding the rule currently in force — and those representations work by biasing the rest of the brain, the same way attention tips a competition among neural representations toward one interpretation. In the Stroop task, where naming the ink color of the word “red” printed in blue pits a practiced habit against an instruction, there is no homunculus reading the word and overriding it. There is a goal, held active, lending its weight to the weaker pathway until it wins.

An AI system built on this model does not need a dedicated executive module enumerating working memory, cognitive flexibility, goal-directed behavior, evaluation and impulse control as five capacities to install. It needs a goal representation with the standing to bias competing processes, which is a narrower and more tractable engineering target. It is a claim about what not to build. It is also the mechanism chapter~\ref{sec:3}'s floor would have to be, if the floor is a commitment rather than a component: not a rule module consulted before acting, but a goal representation with enough standing to bias every competing process toward it. That is why the floor's durability turns out to be a question about disposition. A bias of this kind has no location to attack, and no location to inspect either.

\runin{One mechanism, several jobs}
Plasticity supplies one principle from this literature specific enough to build with, homeostatic plasticity, which holds overall activity in a stable range and is one of the mechanisms proposed against the catastrophic forgetting an unconstrained learner is prone to; the machine side of that problem is worked under continual learning and curriculum learning.

None of these mechanisms is sufficient alone, and integration is not a further module bolted onto the others. What it requires is the same predictive, error-correcting loop implemented consistently across perception, memory, attention and control, so that a corrected prediction about the world propagates into an updated memory, a redirected attention weight, and a revised goal-bias without a separate synchronization mechanism between separately built parts. That is the architectural bet: one mechanism recruited for several jobs, which is more falsifiable than a hand-designed module for every named cognitive process.
